\begin{enumerate}[label=\thesubsection.\arabic*,ref=\thesubsection.\theenumi]
	\item (Bernoulli Distribution)
		For $X \in \cbrak{0,1}$, the pmf is defined as
        \begin{align}
            P_{X}(k)=
                \begin{cases}
                    p &  k=1 \\
                    q & k=0\\
		    0 & \text{otherwise}
                \end{cases}
        \end{align}
\item The moment generating function 
\begin{align}
\begin{split}
	M_{X}(z)\triangleq \mean{z^{-X}} &= p_{X}(1)z^{-1} +p_{X}(0)z^{0}\\
    				&=pz^{-1}+q
\end{split}
\label{eq:EE26BTECH11213-mgf-def}
        \end{align}
	upon substituting from \eqref{eq:EE26BTECH11213-z_transform_def}.
\item 
(Binomial Distribution)	If $X_i$ are independent,  for
\begin{align}
\begin{split}
%\label{eq:EE26BTECH11213-mgf-sum}
	Y \triangleq \sum_{i=1}^{n}X_i, \quad 
	M_{Y}(z)&=\prod_{i=1}^{n}M_{X_i}(z)        \\
	&=\brak{pZ^{-1}+q}^n
\end{split}
\label{eq:EE26BTECH11213-mgf-prod}
\end{align}
\solution 
From \eqref{eq:EE26BTECH11213-mgf-def},
%and
%\eqref{eq:EE26BTECH11213-mgf-sum}
\begin{align}
	M_{Y}(z)=\mean{z^{-\sum_{i=1}^{n}X_i}}
	&=\mean{\prod_{i=1}^{n}z^{-X_{i}}}
	\\
	&=\prod_{i=1}^{n}\mean{z^{-X_{i}}}
\end{align}
yielding \eqref{eq:EE26BTECH11213-mgf-prod}.
\item From  
\eqref{eq:EE26BTECH11213-mgf-prod},
\begin{align}
	\brak{pz^{-1}+q}^n &= \comb{n}{k}p^kq^{n-k}z^{-k}
	\\
         \implies    p_Y\brak{k}&=
	 \begin{cases}
		 \comb{n}{k}p^kq^{n-k} & 0 \le k \le n \\
		 0 & \text{otherwise}
 \end{cases}
\label{eq:EE26BTECH11213-pmf-bern}
\end{align}
%
using the definition in \eqref{eq:EE26BTECH11213-mgf-def} along with
\eqref{eq:EE26BTECH11213-z_transform_def}.
\item
\begin{align}
\begin{split}
\sinh\brak{at}u(t)\system{L}\frac{a}{s^2-a^2}
\\
\cosh\brak{at}u(t)\system{L}\frac{s}{s^2-a^2}
\end{split}
\label{eq:laplaceh}
\end{align}
\solution
From 
\eqref{eq:EE26BTECH11202-Laplace_5},
\begin{align}
	\cos(\j at) u(t)\system{L} \frac{s}{s^2 + \brak{\j a}^2} 
\end{align}
yielding
\eqref{eq:laplaceh}.
%
\item (Geometric Progression)
\begin{align}
	\sum_{i=0}^{n-1}q^i=\frac{q^n-1}{q-1}
	\label{eq:sumofgp}
\end{align}
\item
\begin{align}
\begin{split}
	\sin\brak{x}&=\sum_{i=1}^{n}-1^{i+1}\frac{x^{2i-1}}{\brak{2i-1}!}\\
	\cos\brak{x}&=\sum_{i=1}^{n}-1^{i+1}\frac{x^{i-1}}{{i-1}!}
\end{split}
\label{eq:expansions}
\end{align}
\item If 
\begin{align}
	f\brak{\alpha x,\alpha y}= \alpha^n
	\alpha^nf\brak{x,y},
\end{align}
$f(x,y)$ is homogeneous of degree $n$.  Then,
\begin{align}
    x\frac{\partial f}{\partial x}+y\frac{\partial f}{\partial y}&=nf
    \label{eq:eulerhomo}
\end{align}
\item (fourth order Runge Kutta method)
	Let a function be defined as follows:\\
	$\frac{dy}{dx}=f\brak{x,y}$ , $y\brak{x_0}=y_0.$\\
	the difference equation is given by 
	\begin{align}
		y_{n+1}&=y_n+\frac{h}{6}\brak{k_1+2k_2+2k_3+k_4}\label{eq-rk4}\\
		x_{n+1}&=x_n+h
	\end{align}
	We find $k_1,k_2,k_3,k_4$ by the following formulae:
	\begin{align}
		k_1&=f\brak{x_n,y_n}\\
		k_2&=f\brak{x_n+\frac{h}{2},y_n+k_1\frac{h}{2}}\\
		k_3&=f\brak{x_n=\frac{h}{2},y_n+k_2\frac{h}{2}}\\
		k_4&=f\brak{x_n+h,y_n+hk_3}
	\end{align}	 
\end{enumerate}


